% AUTHOR: Amlal El Mahrouss
% PURPOSE: WG05: The Execution Semantics: On Axioms, Domains, and Authority.

\documentclass[11pt,a4paper]{article}

\usepackage[utf8]{inputenc}
\usepackage[T1]{fontenc}
\usepackage{amsmath,amssymb,amsthm}
\usepackage{hyperref}
\usepackage[margin=1in]{geometry}

\title{The Execution Semantics: On Axioms, Domains, and Authority.}
\author{Amlal El Mahrouss\\
\texttt{amlal@nekernel.org}}
\date{January 2026}

\begin{document}

\bf \maketitle

\begin{center}
	\rule[0.01cm]{17cm}{0.01cm}
\end{center}

\begin{abstract}
This paper presents a foundational framework for execution semantics, consisting of three interconnected theories: the General Harvard Separation Axiom, Execution Domains Theory, and Execution Authority Theory. Together, these establish execution as a primitive concept that cannot be derived from computational models, define boundaries for execution semantics and resource visibility, and formalize the authority governing execution contexts.
\end{abstract}

\begin{center}
	\rule[1cm]{17cm}{0.01cm}
\end{center}

\section{The General Harvard Separation Axiom}

Let $G$ be a theory that formally defines execution semantics (domains, contexts, authority). Let $O$ be any theory of computational behavior.

\subsection{Properties}

\begin{itemize}
    \item If $O$ models computation, then $O$ requires execution to occur.
    \item Therefore $O$ implicitly depends on $G$'s primitives (execution must be defined for computation to be theorized).
    \item $G$ does not depend on $O$ (execution semantics are primitive, not derived from computational models).
\end{itemize}

\subsection{Conclusion}

We cannot derive $G$ from $O$ without circularity:
\begin{itemize}
    \item Deriving execution semantics from $O$ would mean ``execution depends on a theory that assumes execution exists.''
    \item This creates an impossible circular dependency.
    \item Therefore $G$ must be \textbf{axiomatic}—a foundational primitive that cannot be reduced to other computational theories.
\end{itemize}

Any attempt to make $G = O$ or $G \subseteq O$ fails because $O$ already assumes $G$'s primitives exist.

\section{Execution Domains Theory}

\subsection{Abstract}

An execution domain defines a boundary for execution semantics, resource visibility, and control flow.

Let $C$ denote an execution domain in an execution context $E$. Let $D$ be of the same type as $C$ but of a different execution context.

\subsection{Properties}

\begin{itemize}
    \item $C$ shall not be equal to $D$, as $C$ has a different execution context than $D$.
    \item $C$ may be composed of sub-programs within the execution context $E$.
\end{itemize}

\subsection{On Execution Contexts}

Execution contexts are treated as abstract semantic parameters, and execution domains as abstract structures indexed by those parameters.

\section{Execution Authority Theory}

\subsection{Abstract}

An execution domain is defined as previously stated in Section~2. An execution authority is responsible for defining whose semantics may be used for an execution context.

A \emph{trait} is a set of formal rules defining the semantic concepts of an execution context.

Let $A$ be an execution authority of type $T$, such that $T$ is a trait of an execution context.

\subsection{Properties}

Let $C$ denote an execution domain in an execution context $E$. Let $Z$ denote an execution domain in an execution context $X$.

\begin{itemize}
    \item If $X$ does not equal or is not semantically substitutable with $E$—or vice versa—then $C$ shall not equal $Z$.
    \item If $Z$ or $C$ are defined as a null execution context, then the said context—defined as $N$—is not equal to $\lnot N$.
\end{itemize}

\section{Conclusion}

Let $G$ be a theory that formally defines execution semantics—consisting of three foundational theories—the General Harvard Separation Axiom, Execution Domains Theory, and Execution Authority Theory. Together they establish $G$ as a framework that consists of the analysis of Execution Semantics. Or Execution Theory.

\section{References}

\begin{itemize}
	\item El Mahrouss, A. (2026). Methodology for Freestanding Development. Zenodo. https://doi.org/10.5281/zenodo.18362425
	
	\item El Mahrouss, A. (2026). Methodology for Process and Image Computation. (v1.0.0). Zenodo. https://doi.org/10.5281/zenodo.18362503
\end{itemize}

\end{document}
